% 子文件：3.复数与复变函数.tex
% 对应教材：《复变函数与积分变换》(哈工大数学系编, 第三版, 科学出版社) 第 1 章
% 对应讲义：包益欣《复变函数》第 1 章
% 节次：§1.1 复数运算及几何表示（1.1.1 概念与四则运算 / 1.1.2 几何表示 / 1.1.3 共轭复数
%       / 1.1.4 乘除、乘方与开方 / 1.1.5 复球面与无穷远点）
%       §1.2 复平面上的点集（1.2.1 基本概念 / 1.2.2 区域和曲线）
%       §1.3 复变函数（1.3.1 定义与几何意义 / 1.3.2 极限与连续性）
% 由 1.main.tex 用 \input 引入：只写内容，不写 \documentclass、不重复 \usepackage

\chapter{复数与复变函数}

\section{复数运算及几何表示}

\subsection{复数概念及四则运算}
实系数方程 \(x^{2}+1=0\) 在实数域内无解，为此引进虚单位 \(\mathrm{i}\)，规定
\[
\mathrm{i}^{2}=-1 .
\]
称形如
\[
z=x+\mathrm{i}y\qquad (x,y\in\mathbb{R})
\]
的数为\textbf{复数}，\(x,y\) 分别称为 \(z\) 的\textbf{实部}与\textbf{虚部}，记作
\[
x=\operatorname{Re}z,\qquad y=\operatorname{Im}z .
\]
虚部为零的复数即实数，故 \(\mathbb{R}\subset\mathbb{C}\)；实部为零而虚部不为零的复数称为\textbf{纯虚数}。两个复数相等当且仅当实部与虚部分别相等：
\[
z_{1}=z_{2}\iff \operatorname{Re}z_{1}=\operatorname{Re}z_{2}\ \text{且}\ \operatorname{Im}z_{1}=\operatorname{Im}z_{2}.
\]

设 \(z_{1}=x_{1}+\mathrm{i}y_{1}\)，\(z_{2}=x_{2}+\mathrm{i}y_{2}\)，四则运算定义为
\[
z_{1}\pm z_{2}=(x_{1}\pm x_{2})+\mathrm{i}(y_{1}\pm y_{2}),
\]
\keyformula{z_{1}z_{2}=(x_{1}x_{2}-y_{1}y_{2})+\mathrm{i}\,(x_{1}y_{2}+x_{2}y_{1})}
\keyformula{\frac{z_{1}}{z_{2}}=\frac{x_{1}x_{2}+y_{1}y_{2}}{x_{2}^{2}+y_{2}^{2}}
+\mathrm{i}\,\frac{x_{2}y_{1}-x_{1}y_{2}}{x_{2}^{2}+y_{2}^{2}}\qquad (z_{2}\neq 0)}
其中乘法由分配律与 \(\mathrm{i}^{2}=-1\) 直接得到，除法则是“分子分母同乘分母的共轭”的结果。四则运算的结果仍是复数，且加法与乘法满足交换律、结合律与分配律，故 \(\mathbb{C}\) 构成一个域。

由此立得
\[
\operatorname{Re}(z_{1}\pm z_{2})=\operatorname{Re}z_{1}\pm\operatorname{Re}z_{2},\qquad
\operatorname{Im}(z_{1}\pm z_{2})=\operatorname{Im}z_{1}\pm\operatorname{Im}z_{2}.
\]

\corepoint{复数域\textbf{不是有序域}：两个复数（除非都是实数）“谁大谁小”没有意义，不能写 \(z_{1}<z_{2}\)；能比较大小的只有模（非负实数）。另外，\(z_{1}z_{2}=0\iff z_{1}=0\) 或 \(z_{2}=0\) 这一“无零因子”性质在复数域中依然成立，做题时可直接用。}

\subsection{复数的几何表示}
取平面直角坐标系，让复数 \(z=x+\mathrm{i}y\) 与点 \((x,y)\) 对应，此平面称为\textbf{复平面}（\(z\) 平面），\(x\) 轴、\(y\) 轴分别称为\textbf{实轴}、\textbf{虚轴}。于是“点 \(z\)”与“复数 \(z\)”不加区别；点 \(z\) 与原点 \(O\) 的连线所成向量 \(\overrightarrow{Oz}\) 也与 \(z\) 一一对应（零向量对应 \(z=0\)），故复数又可看作平面向量。

向量 \(\overrightarrow{Oz}\) 的长度称为 \(z\) 的\textbf{模}（绝对值）：
\[
|z|=r=\sqrt{x^{2}+y^{2}}\ \ (\geq 0),
\]
实轴正向转到与 \(\overrightarrow{Oz}\) 方向一致时所成的角 \(\theta\) 称为 \(z\) 的\textbf{辐角}，记 \(\operatorname{Arg}z\)。\(z\neq 0\) 时辐角有无穷多个值，彼此相差 \(2\pi\) 的整数倍；\(z=0\) 的模为零，辐角无意义。满足 \(-\pi<\arg z\leqslant \pi\) 的那一个值称为\textbf{辐角主值}，于是
\[
\operatorname{Arg}z=\arg z+2k\pi\qquad (k=0,\pm 1,\pm 2,\cdots).
\]
主值不能只写 \(\arctan\frac{y}{x}\)，须按象限分情况：

\begin{center}
\begin{tabular}{ll}
\toprule
条件 & \(\arg z\) \\
\midrule
\(x>0\) & \(\arctan\dfrac{y}{x}\) \\
\(x<0,\ y\geqslant 0\) & \(\arctan\dfrac{y}{x}+\pi\) \\
\(x<0,\ y<0\) & \(\arctan\dfrac{y}{x}-\pi\) \\
\(x=0,\ y>0\) & \(\dfrac{\pi}{2}\) \\
\(x=0,\ y<0\) & \(-\dfrac{\pi}{2}\) \\
\bottomrule
\end{tabular}
\end{center}

由直角坐标与极坐标的关系 \(x=r\cos\theta,\ y=r\sin\theta\) 得\textbf{三角表示式}
\[
z=r(\cos\theta+\mathrm{i}\sin\theta),
\]
再由欧拉公式 \(\mathrm{e}^{\mathrm{i}\theta}=\cos\theta+\mathrm{i}\sin\theta\) 得\textbf{指数表示式}
\[
z=r\,\mathrm{e}^{\mathrm{i}\theta}.
\]
两复数相等 \(\iff\) 模相等且辐角相差 \(2\pi\) 的整数倍（主值相等）。

几何上，\(z_{1}+z_{2}\) 与 \(z_{1}-z_{2}\) 分别是平行四边形（以 \(\overrightarrow{Oz_{1}},\overrightarrow{Oz_{2}}\) 为邻边）的两条对角线向量；\(|z_{1}-z_{2}|\) 就是复平面上两点 \(z_{1},z_{2}\) 之间的距离。由此得到两个基本不等式：
\keyformula{\bigl||z_{1}|-|z_{2}|\bigr|\leqslant |z_{1}\pm z_{2}|\leqslant |z_{1}|+|z_{2}|}
同时有 \(|x|\leqslant |z|\)，\(|y|\leqslant |z|\)，\(|z|\leqslant |x|+|y|\)，即
\[
|\operatorname{Re}z|\leqslant |z|,\qquad |\operatorname{Im}z|\leqslant |z|.
\]

\subsection{共轭复数}
实部相同、虚部互为相反数的两个复数称为\textbf{共轭复数}，\(z=x+\mathrm{i}y\) 的共轭记作
\[
\bar z=x-\mathrm{i}y .
\]
共轭运算与四则运算可交换次序：
\[
\overline{z_{1}\pm z_{2}}=\bar z_{1}\pm\bar z_{2},\qquad
\overline{z_{1}z_{2}}=\bar z_{1}\bar z_{2},\qquad
\overline{\left(\frac{z_{1}}{z_{2}}\right)}=\frac{\bar z_{1}}{\bar z_{2}}\quad(z_{2}\neq 0).
\]
并且
\keyformula{z+\bar z=2\operatorname{Re}z,\qquad z-\bar z=2\mathrm{i}\operatorname{Im}z,\qquad z\bar z=|z|^{2}}
由最后一式，除法可写成 \(\dfrac{z_{1}}{z_{2}}=\dfrac{z_{1}\bar z_{2}}{|z_{2}|^{2}}\)，这是实际计算中最常用的形式。又 \(z\) 为实数 \(\iff z=\bar z\)。利用共轭还可证明 \(|z_{1}z_{2}|=|z_{1}||z_{2}|\)、\(|z_{1}/z_{2}|=|z_{1}|/|z_{2}|\)（模是可乘的）。

\subsection{乘除、乘方与开方}
设 \(z_{1}=r_{1}\mathrm{e}^{\mathrm{i}\theta_{1}}\)，\(z_{2}=r_{2}\mathrm{e}^{\mathrm{i}\theta_{2}}\)，则
\keyformula{z_{1}z_{2}=r_{1}r_{2}\,\mathrm{e}^{\mathrm{i}(\theta_{1}+\theta_{2})},\qquad
\frac{z_{1}}{z_{2}}=\frac{r_{1}}{r_{2}}\,\mathrm{e}^{\mathrm{i}(\theta_{1}-\theta_{2})}\quad(z_{2}\neq 0)}
即：\textbf{两复数相乘，模相乘、辐角相加；相除则模相除、辐角相减}。几何意义是“伸缩 \(+\) 旋转”：乘 \(z_{2}\) 把 \(z_{1}\) 的模放大 \(r_{2}\) 倍并逆时针旋转 \(\theta_{2}\)。

反复应用乘法法则即得\textbf{棣莫弗（De Moivre）公式}
\[
z^{n}=r^{n}\mathrm{e}^{\mathrm{i}n\theta}=r^{n}(\cos n\theta+\mathrm{i}\sin n\theta),\qquad n=1,2,\cdots
\]
以及 \((\cos\theta+\mathrm{i}\sin\theta)^{n}=\cos n\theta+\mathrm{i}\sin n\theta\)。

\textbf{开方}是乘方的逆运算：求 \(w\) 使 \(w^{n}=z\)（\(z\neq 0\)）。设 \(w=\rho\mathrm{e}^{\mathrm{i}\varphi}\)，则 \(\rho^{n}=r\)、\(n\varphi=\theta+2k\pi\)，故
\keyformula{w_{k}=\sqrt[n]{r}\;\mathrm{e}^{\mathrm{i}\frac{\theta+2k\pi}{n}},\qquad k=0,1,2,\cdots,n-1}
非零复数开 \(n\) 次方恰有 \(n\) 个互不相同的值，它们的模都等于 \(\sqrt[n]{r}\)，辐角依次相差 \(\frac{2\pi}{n}\)，即均匀分布在以原点为心、半径 \(\sqrt[n]{r}\) 的圆周上（内接正 \(n\) 边形的顶点）。特别地，\(\sqrt[n]{1}\) 的 \(n\) 个值称为 \(n\) 次\textbf{单位根}，即
\[
\omega_{k}=\mathrm{e}^{\mathrm{i}\frac{2k\pi}{n}}\quad (k=0,1,\cdots,n-1),\qquad
1+z+z^{2}+\cdots+z^{n-1}=\frac{z^{n}-1}{z-1}\ (z\neq 1).
\]

\corepoint{\(\operatorname{Arg}(z_{1}z_{2})=\operatorname{Arg}z_{1}+\operatorname{Arg}z_{2}\) 只在“集合意义”下成立（两边都是无穷集合）；对\textbf{主值} \(\arg\) 一般\textbf{不成立}：取 \(z_{1}=z_{2}=-1\)，则 \(\arg(z_{1}z_{2})=\arg 1=0\)，而 \(\arg z_{1}+\arg z_{2}=2\pi\)，两者相差 \(2\pi\)。凡涉及主值的等式，都要先验证是否跨越了主值区间 \((-\pi,\pi]\) 的“割线”。

开方千万别漏根：解 \(z^{n}=a\) 时答案必须是 \(n\) 个值，写成 \(\sqrt[n]{r}\,\mathrm{e}^{\mathrm{i}\theta/n}\) 只给出一个根，这是最常见的失分点。}

\subsection{复球面与无穷远点}
在复平面上方放一个与之相切于原点的球面 \(S\)（如取球面 \(x^{2}+y^{2}+(u-1)^{2}=1\)，北极点 \(N=(0,0,2)\)），作\textbf{球极投影}：复平面上任一点 \(z\) 与 \(N\) 的连线交球面于唯一异于 \(N\) 的点 \(P(z)\)。这样得到 \(\mathbb{C}\) 到 \(S\setminus\{N\}\) 的一一对应，且该对应是连续同胚（逆映射也连续）。

再在 \(\mathbb{C}\) 中添加一个\textbf{无穷远点} \(\infty\)，规定 \(z=\infty\) 对应北极 \(N\)，便得
\[
\mathbb{C}_{\infty}=\mathbb{C}\cup\{\infty\},
\]
它与整个球面 \(S\) 一一对应，故称 \(S\) 为\textbf{复球面}，\(\mathbb{C}_{\infty}\) 为\textbf{扩充复平面}。

\begin{itemize}[leftmargin=2em]
  \item \(z\to\infty\) 指 \(|z|\to+\infty\)，与趋近方向无关；因此扩充复平面中\textbf{只有一个} \(\infty\)，没有 \(+\infty\)、\(-\infty\)、\(\mathrm{i}\infty\) 之分。
  \item \(\infty\) 的 \(\delta\) 邻域规定为 \(\{z: |z|>M\}\cup\{\infty\}\)（\(M=1/\delta>0\)），其去心邻域即 \(\{z: |z|>M\}\)。
  \item 复平面上任一直线（作为点集）都以 \(\infty\) 为聚点；无界区域在球面上以 \(\infty\) 为边界点或内点。
  \item 运算约定：\(\infty\pm z=\infty\)，\(\infty\cdot z=\infty\ (z\neq 0)\)，\(\frac{z}{\infty}=0\)，\(\frac{\infty}{0}=\infty\)；而 \(\infty\pm\infty\)、\(0\cdot\infty\)、\(\frac{0}{0}\)、\(\frac{\infty}{\infty}\) 无意义。
\end{itemize}

\section{复平面上的点集}

复变函数的定义域是复平面（或复球面）上的点集，故先约定点集的术语。

\subsection{基本概念}
设 \(z_{0}\in\mathbb{C}\)，\(\delta>0\)，称
\[
B(z_{0},\delta)=\{z: |z-z_{0}|<\delta\},\qquad
B^{\circ}(z_{0},\delta)=\{z: 0<|z-z_{0}|<\delta\}
\]
为 \(z_{0}\) 的 \(\delta\) \textbf{邻域}与\textbf{去心 \(\delta\) 邻域}。

设 \(E\subset\mathbb{C}\)：
\begin{itemize}[leftmargin=2em]
  \item \(z\) 是 \(E\) 的\textbf{内点}：存在 \(\delta>0\) 使 \(B(z,\delta)\subset E\)；\(E\) 的内点全体称为 \(E\) 的\textbf{内部}。
  \item \(E\) 为\textbf{开集}：\(E\) 的每一点都是内点。
  \item \(z\) 是 \(E\) 的\textbf{边界点}：对任意 \(\delta>0\)，\(B(z,\delta)\) 内既有 \(E\) 的点又有不属于 \(E\) 的点；边界点全体记 \(\partial E\)。
  \item \(E\) 的\textbf{闭包} \(\overline{E}=E\cup\partial E\)；若 \(\overline{E}=E\)（即 \(E\) 含其全部边界点）则 \(E\) 为\textbf{闭集}。
  \item \(z\) 是 \(E\) 的\textbf{聚点}（极限点）：对任意 \(\delta>0\)，\(B^{\circ}(z,\delta)\cap E\neq\varnothing\)。
  \item \(z\) 是 \(E\) 的\textbf{孤立点}：\(z\in E\) 且存在 \(\delta>0\) 使 \(B^{\circ}(z,\delta)\cap E=\varnothing\)。
  \item \(E\) \textbf{有界}：存在 \(M>0\) 使 \(|z|\leqslant M\) 对一切 \(z\in E\) 成立。
\end{itemize}

\begin{center}
\begin{tabular}{ll}
\toprule
命题 & 内容 \\
\midrule
内点与聚点 & \(E\) 的内点必是 \(E\) 的聚点 \\
孤立点与边界 & \(E\) 的孤立点必是 \(E\) 的边界点，且不是聚点 \\
闭集判别 & \(E\) 闭 \(\iff E\) 含其一切聚点 \\
\bottomrule
\end{tabular}
\end{center}

点 \(z_{0}\) 到集合 \(E\) 的\textbf{距离}定义为 \(\operatorname{dist}(z_{0},E)=\inf\{|z-z_{0}|: z\in E\}\)。

\subsection{区域和曲线}
\textbf{区域}（开区域）指\textbf{连通的开集}，其中“连通”指集合内任意两点可用集合内的折线连接。区域 \(D\) 连同其边界 \(\partial D\) 称为\textbf{闭区域} \(\overline{D}=D\cup\partial D\)。含于某个圆内的集合为有界集，相应的区域称为有界区域，否则为无界区域。

\textbf{曲线}是闭区间 \([a,b]\) 到复平面的连续映射
\[
C:\ z=z(t)=x(t)+\mathrm{i}y(t),\qquad a\leqslant t\leqslant b,
\]
\(z(a)\)、\(z(b)\) 称为起点与终点。若 \(z(t)\) 在 \([a,b]\) 上除端点外处处不同，称为\textbf{简单曲线}（无重点）；若还有 \(z(a)=z(b)\)，称为\textbf{简单闭曲线}。若 \(z'(t)\) 在 \([a,b]\) 上连续且处处不为零，称 \(C\) 为\textbf{光滑曲线}；有限条光滑曲线首尾相接称为\textbf{逐段光滑曲线}。

\textbf{若尔当（Jordan）定理}\quad 一条简单闭曲线把复平面分成内部（有界区域）与外部（无界区域）两部分，曲线本身是二者的公共边界。

\textbf{单连通区域}：区域 \(D\) 内任一条简单闭曲线的内部都含于 \(D\)。圆盘、半平面、带形区域、角形区域都是单连通区域；去心圆盘、圆环是多连通区域（其边界由两条以上互不相交的闭曲线组成）。

对区域的边界定向：沿边界前进时\textbf{区域始终在左侧}的方向为\textbf{正向}。于是单连通区域边界逆时针为正；圆环区域的外边界逆时针为正、内边界顺时针为正。

\corepoint{“区域”必须同时满足\textbf{开}与\textbf{连通}，缺一不可：\(\{z: |z|\leqslant 1\}\) 不开，两个不相交圆盘的并集不连通，都不是区域。求解析区域、判断积分能否用柯西定理时，第一步永远是看清定义域是开区域还是闭区域、单连通还是多连通——多连通区域上“绕洞”的积分一般不为零。}

\section{复变函数}

\subsection{定义与几何意义}
设 \(E\subset\mathbb{C}\)。若对每个 \(z\in E\)，按某一法则 \(f\) 有确定的复数 \(w\) 与之对应，则称 \(w=f(z)\) 是定义在 \(E\) 上的（单值）\textbf{复变函数}，\(E\) 为定义域，\(\{f(z): z\in E\}\) 为值域。若对每个 \(z\in E\) 对应的是一个复数集合 \(f(z)\)，则称 \(f\) 为\textbf{多值函数}。多值函数 \(f\) 在 \(E\) 上的\textbf{单值分支}指这样的单值函数 \(f_{0}\)：对每个 \(z\in E\) 都有 \(f_{0}(z)\in f(z)\)。

记 \(z=x+\mathrm{i}y\)，\(w=u+\mathrm{i}v\)，则
\[
w=f(z)=u(x,y)+\mathrm{i}v(x,y),
\]
即一个复变函数等价于两个实值二元函数 \(u,v\)。

\textbf{几何意义}：\(w=f(z)\) 把 \(z\) 平面上的点集 \(E\) 映为 \(w\) 平面上的点集 \(f(E)\)，称为\textbf{映射}（变换）。若 \(f\) 是 \(E\) 到 \(f(E)\) 的一一对应，则存在\textbf{反函数} \(z=f^{-1}(w)\)；两个映射可复合。研究映射 \(w=f(z)\) 时，通常同时画出 \(z\) 平面与 \(w\) 平面，看“把什么图形映成什么图形”——这正是第 6 章保形映射的出发点。

\subsection{极限与连续性}
设 \(f\) 在 \(z_{0}\) 的去心邻域内有定义。若存在复数 \(A\)，使对任意 \(\varepsilon>0\)，存在 \(\delta>0\)，当 \(0<|z-z_{0}|<\delta\) 时恒有
\[
|f(z)-A|<\varepsilon,
\]
则称 \(z\to z_{0}\) 时 \(f(z)\) 以 \(A\) 为\textbf{极限}，记 \(\lim\limits_{z\to z_{0}}f(z)=A\)。这里 \(z_{0}\) 与 \(A\) 都可以是 \(\infty\)（此时把相应的绝对值不等式换成 \(|z|>M\) 或 \(|f(z)|>M\)）。

极限的四则运算、唯一性与实函数相同；关键差别在于 \(z\to z_{0}\) 是\textbf{在平面上沿任意路径}趋近，故有\textbf{归结原则}：\(\lim\limits_{z\to z_{0}}f(z)=A\) 的充要条件是对任一趋于 \(z_{0}\) 的点列 \(\{z_{n}\}\) 都有 \(f(z_{n})\to A\)。

由于 \(f=u+\mathrm{i}v\)，极限问题可化为两个实二元函数的极限：
\keyformula{\lim_{z\to z_{0}}f(z)=A+\mathrm{i}B\iff
\lim_{(x,y)\to(x_{0},y_{0})}u(x,y)=A\ \text{且}\ \lim_{(x,y)\to(x_{0},y_{0})}v(x,y)=B}
若 \(\lim\limits_{z\to z_{0}}f(z)=f(z_{0})\)（要求 \(f\) 在 \(z_{0}\) 的某邻域内有定义），则称 \(f\) 在 \(z_{0}\) 处\textbf{连续}；在区域 \(D\) 内每点连续则称 \(f\) 在 \(D\) 上连续。

由复极限与实极限的等价性立得：
\begin{itemize}[leftmargin=2em]
  \item \(f\) 在 \(z_{0}\) 连续 \(\iff u,v\) 在 \((x_{0},y_{0})\) 连续；
  \item 连续函数的和、差、积、商（分母不为零处）与复合仍连续；
  \item 在有界闭集（紧集）上连续的复函数有界，且 \(|f(z)|\) 能取到最大值；并且一致连续。
\end{itemize}

\corepoint{判断“极限不存在”的标准手段是\textbf{沿特殊路径}：只要找到两条不同的路径使 \(f(z)\) 趋于不同的值（或沿某条路径无极限），极限即不存在。典型反例：
\[
f(z)=\frac{\bar z}{z}\ (z\neq 0),
\]
沿实轴 \(z=x\to 0\) 得 \(f\to 1\)，沿虚轴 \(z=\mathrm{i}y\to 0\) 得 \(f\to -1\)，故 \(z=0\) 处极限不存在。

另一个常考陷阱是\textbf{辐角主值不连续}：\(\arg z\) 在负实轴上的点（含 \(z<0\) 的实轴）不连续（上半平面趋近得 \(\pi\)，下半平面趋近得 \(-\pi\)），但在割开的区域 \(\mathbb{C}\setminus(-\infty,0]\) 上连续。它直接影响 \(w=\ln z\)、\(w=z^{a}\) 的解析区域，第 2 章会反复用到。}
